\chapter{Warstwa sprzętowa stacji bazowej}
\label{chap:sprzet}


\section{Dobór podzespołów}
\label{sec:dobor}
% ANALIZA porównawcza na podstawie danych katalogowych

Dobór podzespołów przeprowadzono na podstawie analizy porównawczej parametrów
katalogowych dostępnych rozwiązań, przyjmując kryteria wynikające z~wymagań
przedstawionych w~podrozdziale~\ref{sec:wymagania}. Jako dodatkowe kryterium
uwzględniono dostępność podzespołów, co pozwoliło skrócić czas budowy
prototypu i~ograniczyć jego koszt. Poprawność dokonanego wyboru zweryfikowano
pomiarami przedstawionymi w~rozdziale~\ref{chap:badania}.

\begin{table}[H]
  \centering
  \caption{Porównanie platform obliczeniowych}
  \label{tab:platformy}
  \small
  \begin{tabular}{l c c c c c}
    \toprule
    Platforma & Pobór mocy [W] & GPIO/UART/SPI & ROS~2 Jazzy & GPU & Koszt \\
    \midrule
    Raspberry Pi 4B       & -- & -- & -- & -- & -- \\
    Raspberry Pi 5        & -- & -- & -- & -- & -- \\
    NVIDIA Jetson Orin Nano & -- & -- & -- & -- & -- \\
    Mini-PC x86           & -- & -- & -- & -- & -- \\
    \bottomrule
  \end{tabular}
\end{table}

\begin{table}[H]
  \centering
  \caption{Porównanie odbiorników SDR}
  \label{tab:sdr}
  \small
  \begin{tabular}{l c c c c c}
    \toprule
    Odbiornik & Zakres [MHz] & ADC [bit] & Filtr 1090~MHz & dump1090 & Koszt \\
    \midrule
    RTL-SDR Blog V3             & -- & -- & -- & -- & -- \\
    FlightAware Pro Stick Plus  & -- & -- & -- & -- & -- \\
    Airspy Mini                 & -- & -- & -- & -- & -- \\
    HackRF One                  & -- & -- & -- & -- & -- \\
    \bottomrule
  \end{tabular}
\end{table}

Wybrane podzespoły zestawiono w~tabeli~\ref{tab:podzespoly}.

\begin{table}[H]
  \centering
  \caption{Zestawienie podzespołów stacji bazowej}
  \label{tab:podzespoly}
  \small
  \begin{tabular}{l l >{\raggedright\arraybackslash}p{5.5cm}}
    \toprule
    Podzespół & Model & Uwagi \\
    \midrule
    Komputer        & Raspberry Pi 4B          & system uruchamiany z~pamięci USB \SI{64}{GB} \\
    Odbiornik SDR   & RTL-SDR Blog V3          & złącze USB \\
    Antena          & dipol z~zestawu RTL-SDR Blog & ramiona ok.\ \SI{6.5}{cm}, montaż pionowy \\
    Odbiornik GPS   & moduł z~wyjściem NMEA    & UART, \SI{9600}{bit/s}, komunikaty GGA i~RMC co \SI{1}{s} \\
    Wyświetlacz     & TFT ILI9341, $240\times320$ & przekątna 2,4~cala, magistrala SPI \\
    Zasilanie       & powerbank LYNXER RPP-272 & wyjście USB \SI{5}{V} \\
    \bottomrule
  \end{tabular}
\end{table}

% ---------------------------------------------------------------------
\section{Tor odbioru radiowego}
\label{sec:tor_radiowy}
% Antena dipolowa ćwierćfalowa (odwołanie do wzoru
% \eqref{eq:lambda}), polaryzacja pionowa, RTL-SDR, zakłócenia od USB,
% umiejscowienie anteny.

% ---------------------------------------------------------------------
\section{Odbiornik GPS}
\label{sec:odbiornik_gps}
%Połączenie UART, ustawienia, komunikaty, antena z widokiem
% na niebo, utrata fixa w pomieszczeniu.

% ---------------------------------------------------------------------
\section{Elementy obsługi i~schemat połączeń}
\label{sec:polaczenia}
% Wyświetlacz, dioda, przycisk; schemat (rysunek);
% konflikt GPIO3 z I2C (dlaczego I2C wyłączone).

Połączenia elementów z~wyprowadzeniami złącza GPIO komputera Raspberry Pi
zestawiono w~tabeli~\ref{tab:gpio}.

\begin{table}[H]
  \centering
  \caption{Połączenia ze złączem GPIO komputera Raspberry Pi}
  \label{tab:gpio}
  \small
  \begin{tabular}{>{\raggedright}p{3.6cm} >{\raggedright}p{5.2cm} >{\raggedright\arraybackslash}p{4.8cm}}
    \toprule
    Element & Wyprowadzenia (numer fizyczny) & Uwagi \\
    \midrule
    Dioda stanu          & 11 (GPIO17), 34 (GND) & rezystor szeregowy \SIrange{220}{330}{\ohm} \\
    Przycisk zasilania   & 5 (GPIO3), 6 (GND)    & wybudzanie ze stanu zatrzymania \\
    Wyświetlacz (SPI0)   & 19, 21, 23, 24; DC: 18 (GPIO24); RST: 22 (GPIO25) & sterownik ILI9341 \\
    Odbiornik GPS (UART) & 8, 10                 & \texttt{/dev/ttyAMA0} \\
    \bottomrule
  \end{tabular}
\end{table}

% ---------------------------------------------------------------------
\section{Zasilanie i~obudowa}
\label{sec:zasilanie}
% Powerbank, obudowa drukowana 3D, chłodzenie, okienko na
% wskaźnik powerbanka, dostęp do ładowania.